\documentclass[tikz,border=20pt]{standalone}
\usepackage[utf8]{inputenc}
\usetikzlibrary{shapes.geometric, arrows.meta, positioning, shadows, calc, fit, backgrounds, decorations.pathreplacing}

% Define Professional Colors with better contrast
\definecolor{layerCognitive}{HTML}{2C3E50}  % Dark Blue/Grey
\definecolor{layerNetwork}{HTML}{F39C12}    % Orange
\definecolor{layerPhysics}{HTML}{C0392B}    % Deep Red
\definecolor{layerOutput}{HTML}{27AE60}     % Green
\definecolor{boxFill}{HTML}{FDFFE6}         % Very light cream for boxes
\definecolor{lineColor}{HTML}{34495E}

\begin{document}

\begin{tikzpicture}[
    node distance=1.5cm and 1.0cm,
    font=\sffamily,
    % Core Styles
    block/.style={
        rectangle, 
        rounded corners=4pt, 
        fill=boxFill,
        draw=lineColor,
        very thick,
        text width=3.2cm, 
        minimum height=1.6cm, 
        align=center, 
        drop shadow={opacity=0.3},
        font=\small
    },
    header/.style={
        rectangle,
        rounded corners=2pt,
        text=white,
        font=\bfseries\Large,
        align=center,
        minimum height=1.0cm,
        minimum width=16cm,
        inner sep=5pt
    },
    % Connectors
    arrow/.style={very thick, ->, >=Stealth, color=lineColor},
    feedback/.style={very thick, dashed, ->, >=Stealth, color=layerPhysics}
]

% =================================================================================
% LAYER 1: COGNITIVE AGENTS
% =================================================================================

\node[header, fill=layerCognitive] (Header1) {LAYER 1: COGNITIVE AGENTS (User Behavior)};

% Personas (Tallest element determining layer height)
\node[block, below=1.0cm of Header1.west, anchor=north west, xshift=0.5cm, text width=3.5cm] (Personas) {
    \textbf{5 Personas}\\
    \raggedright
    \scriptsize
    1. Gamer Gary\\
    2. Creator Chloe\\
    3. Commuter Chris\\
    4. Chatter Emma\\
    5. Streamer Steve
};

% LLM Agent
\node[block, right=1.0cm of Personas, text width=4.5cm] (LLM) {
    \textbf{Learning LLM Agent}\\
    \scriptsize (GPT-4 + Memory)\\
    \vspace{0.1cm}
    \textit{"Reflects \& Adapts Strategies based on Thermal History"}
};

% RL Agent
\node[block, right=1.0cm of LLM, text width=4.5cm] (RL) {
    \textbf{Deep RL Agent (DQN)}\\
    \scriptsize Double DQN + Epsilon Greedy\\
    \vspace{0.1cm}
    \textit{"Optimizes Satisfaction vs Power Consumption"}
};

% Internal Connections
\draw[arrow] (Personas) -- (LLM);
\draw[arrow] (LLM) -- node[above, font=\scriptsize] {Guidance} (RL);

% Draw Frame for Layer 1
\begin{scope}[on background layer]
    \node[draw=layerCognitive, dashed, very thick, inner sep=10pt, fit=(Header1) (Personas) (RL)] (Layer1Frame) {};
\end{scope}

% =================================================================================
% LAYER 2: NETWORK & HARDWARE
% =================================================================================

% Important: Position Header 2 relative to the FRAME of Layer 1, not the Header 1
\node[header, fill=layerNetwork, below=1.5cm of Layer1Frame.south] (Header2) {LAYER 2: NETWORK \& HARDWARE LOADING};

% RRC State Machine
\node[block, below=1.0cm of Header2.west, anchor=north west, xshift=1.5cm] (RRC) {
    \textbf{5G RRC State Machine}\\
    \scriptsize IDLE $\leftrightarrow$ CONNECTED $\leftrightarrow$ TAIL\\
    \vspace{0.1cm}
    \textit{"The Tail Energy Penalty"}
};

% GPU Model
\node[block, right=1.0cm of RRC] (GPU) {
    \textbf{APGPM GPU Model}\\
    \scriptsize Render / Compute / AI Loads\\
    \vspace{0.1cm}
    \textit{Thermal Throttling Logic}
};

% Fast Charging
\node[block, right=1.0cm of GPU] (Charging) {
    \textbf{Fast Charging Tracker}\\
    \scriptsize Stress Factors $\sigma(I)$\\
    \vspace{0.1cm}
    \textit{LifecyclE 2025 Model}
};

% Connections Layer 1 -> Layer 2 (Action Bus)
% Route from RL Agent -> Right -> Down (Past Header 2) -> Left -> Down
\draw[arrow] (RL.east) -- ++(1.0,0) coordinate(busStart) 
    -- (busStart |- Header2.south) coordinate(busTurn) % Go down to below Header 2
    -- ++(0,-0.5) coordinate(busUnder) % Go a bit lower
    -| (RRC.north); % Connect to RRC

\draw[arrow] (busUnder) -| (GPU.north);
\draw[arrow] (busUnder) -| (Charging.north);

% Label for the Action
\node[right, font=\bfseries\small, color=lineColor] at ($(busStart)!0.5!(busTurn)$) {Action $A_t$};

% Layer 2 Frame
\begin{scope}[on background layer]
    \node[draw=layerNetwork, dashed, very thick, inner sep=10pt, fit=(Header2) (RRC) (Charging)] (Layer2Frame) {};
\end{scope}

% =================================================================================
% LAYER 3: UNIFIED PHYSICS ENGINE
% =================================================================================

\node[header, fill=layerPhysics, below=1.5cm of Layer2Frame.south] (Header3) {LAYER 3: UNIFIED PHYSICS ENGINE (ODE Solver)};

% KiBaM
\node[block, below=1.0cm of Header3.west, anchor=north west, xshift=1.5cm] (KiBaM) {
    \textbf{Model I: KiBaM}\\
    \scriptsize Two-Tank Recovery\\
    \vspace{0.1cm}
    $k(h_1 - h_2)$
};

% ECM
\node[block, right=1.0cm of KiBaM] (ECM) {
    \textbf{Model II: ET-SSF}\\
    \scriptsize 2-RC Pairs + OCV(SOC)\\
    \vspace{0.1cm}
    $V_{term} = V_{oc} - IR - V_{p}$
};

% Thermal
\node[block, right=1.0cm of ECM] (Thermal) {
    \textbf{Thermal Dynamics}\\
    \scriptsize Arrhenius Coupling\\
    \vspace{0.1cm}
    $R(T) = R_{ref} \exp(E_a/kT)$
};

% Connections Layer 2 -> Layer 3
\draw[arrow] (RRC.south) -- node[right, font=\tiny] {Power} (ECM.north);
\draw[arrow] (GPU.south) -- node[right, font=\tiny] {Heat} (Thermal.north);
\draw[arrow] (ECM) -- node[above, font=\tiny] {Joule Heat} (Thermal);
\draw[arrow] (ECM) -- node[above, font=\tiny] {I(t)} (KiBaM);

% Layer 3 Frame
\begin{scope}[on background layer]
    \node[draw=layerPhysics, dashed, very thick, inner sep=10pt, fit=(Header3) (KiBaM) (Thermal)] (Layer3Frame) {};
\end{scope}

% =================================================================================
% LAYER 4: OUTPUTS
% =================================================================================

\node[header, fill=layerOutput, below=1.5cm of Layer3Frame.south] (Header4) {LAYER 4: OUTPUTS & STRATEGIES};

\node[block, below=0.5cm of Header4, text width=14cm, minimum height=2cm] (Results) {
    \textbf{Key Insights \& Optimizations}\\
    \raggedright
    \begin{itemize}
        \item \textbf{Smart-5G Throttling:} Shift background sync to LTE reduces tail energy (-15\%).
        \item \textbf{Batching:} Grouping notifications reduces RRC state transitions.
        \item \textbf{Emergent Behavior:} LLM agent learns to pre-cool phone before gaming.
    \end{itemize}
};

% Layer 4 Frame
\begin{scope}[on background layer]
    \node[draw=layerOutput, dashed, very thick, inner sep=10pt, fit=(Header4) (Results)] (Layer4Frame) {};
\end{scope}


% =================================================================================
% GLOBAL FEEDBACK LOOP
% =================================================================================

% =================================================================================
% GLOBAL FEEDBACK LOOP
% =================================================================================

% Route: ECM.south -> Down (below Layer 3 frame) -> Right (margin) -> Up -> RL.north
\draw[feedback] (ECM.south) 
    -- ++(0,-0.8) coordinate(fbLow) % Drop below Layer 3 blocks
    -| ($(Layer1Frame.east) + (1.5,0)$) coordinate(fbTurn) % Go Right to margin then Up
    |- (RL.north); % Connect to top of RL Agent

% Label on the vertical rising segment
\node[above, rotate=90, font=\bfseries\large, color=layerPhysics] at ($(fbTurn)!0.7!(fbTurn |- RL.north)$) {State Feedback: $S_{t+1}$ (SOC, Temp)};

\end{tikzpicture}
\end{document}
